% Options for packages loaded elsewhere
\PassOptionsToPackage{unicode}{hyperref}
\PassOptionsToPackage{hyphens}{url}
\PassOptionsToPackage{dvipsnames,svgnames,x11names}{xcolor}
\documentclass[
  11pt,
  a4paper,
]{article}
\usepackage{xcolor}
\usepackage[landscape,margin=18mm,headheight=15pt]{geometry}
\usepackage{amsmath,amssymb}
\setcounter{secnumdepth}{-\maxdimen} % remove section numbering
\usepackage{iftex}
\ifPDFTeX
  \usepackage[T1]{fontenc}
  \usepackage[utf8]{inputenc}
  \usepackage{textcomp} % provide euro and other symbols
\else % if luatex or xetex
  \usepackage{unicode-math} % this also loads fontspec
  \defaultfontfeatures{Scale=MatchLowercase}
  \defaultfontfeatures[\rmfamily]{Ligatures=TeX,Scale=1}
\fi
\usepackage{lmodern}
\ifPDFTeX\else
  % xetex/luatex font selection
\fi
% Use upquote if available, for straight quotes in verbatim environments
\IfFileExists{upquote.sty}{\usepackage{upquote}}{}
\IfFileExists{microtype.sty}{% use microtype if available
  \usepackage[]{microtype}
  \UseMicrotypeSet[protrusion]{basicmath} % disable protrusion for tt fonts
}{}
\makeatletter
\@ifundefined{KOMAClassName}{% if non-KOMA class
  \IfFileExists{parskip.sty}{%
    \usepackage{parskip}
  }{% else
    \setlength{\parindent}{0pt}
    \setlength{\parskip}{6pt plus 2pt minus 1pt}}
}{% if KOMA class
  \KOMAoptions{parskip=half}}
\makeatother
\usepackage{longtable,booktabs,array}
\usepackage{caption}
\captionsetup[table]{skip=6pt}
\newcounter{none} % for unnumbered tables
\usepackage{calc} % for calculating minipage widths
% Correct order of tables after \paragraph or \subparagraph
\usepackage{etoolbox}
\makeatletter
\patchcmd\longtable{\par}{\if@noskipsec\mbox{}\fi\par}{}{}
\makeatother
% Allow footnotes in longtable head/foot
\IfFileExists{footnotehyper.sty}{\usepackage{footnotehyper}}{\usepackage{footnote}}
\makesavenoteenv{longtable}
\setlength{\emergencystretch}{3em} % prevent overfull lines
\providecommand{\tightlist}{%
  \setlength{\itemsep}{0pt}\setlength{\parskip}{0pt}}
\usepackage{fontspec}
\usepackage{fancyhdr}
\usepackage{enumitem}
\usepackage{etoolbox}
\usepackage{xurl}
\usepackage{needspace}
\setmainfont[Path=/usr/local/texlive/2026/texmf-dist/fonts/opentype/public/tex-gyre/,UprightFont=texgyrepagella-regular.otf,BoldFont=texgyrepagella-bold.otf,ItalicFont=texgyrepagella-italic.otf,BoldItalicFont=texgyrepagella-bolditalic.otf]{texgyrepagella-regular.otf}
\setsansfont[Path=/usr/local/texlive/2026/texmf-dist/fonts/opentype/public/tex-gyre/,UprightFont=texgyreheros-regular.otf,BoldFont=texgyreheros-bold.otf,ItalicFont=texgyreheros-italic.otf,BoldItalicFont=texgyreheros-bolditalic.otf]{texgyreheros-regular.otf}
\setmonofont[Path=/usr/local/texlive/2026/texmf-dist/fonts/opentype/SIL/jetbrainsmono-otf/,UprightFont=JetBrainsMono-Regular.otf,BoldFont=JetBrainsMono-Bold.otf,ItalicFont=JetBrainsMono-Italic.otf,BoldItalicFont=JetBrainsMono-BoldItalic.otf,Scale=0.85]{JetBrainsMono-Regular.otf}
\definecolor{notesblue}{HTML}{174A7E}
\definecolor{notesgray}{HTML}{5D6773}
\definecolor{noteslink}{HTML}{1769AA}
\AtBeginDocument{\hypersetup{linkcolor=noteslink,urlcolor=noteslink,pdftitle={VAEs for Video Generation Models},pdfsubject={Compression, architecture, causality, and chunked decoding}}}
\setlength{\parindent}{0pt}
\setlength{\parskip}{0.5em}
\linespread{1.06}
\setcounter{secnumdepth}{0}
\setlist{leftmargin=2em,itemsep=0.2em}
\pagestyle{fancy}
\fancyhf{}
\fancyhead[L]{\small\sffamily\color{notesgray}Video VAE / Architecture Notes}
\fancyhead[R]{\small\sffamily\color{notesgray}Purshow\textquoteright{}s Notes}
\fancyfoot[R]{\small\thepage}
\renewcommand{\headrulewidth}{0.3pt}
\newcommand{\NotesTitle}[1]{{\sffamily\bfseries\color{notesblue}\fontsize{25}{32}\selectfont #1\par}\vspace{2mm}{\color{notesblue}\rule{\linewidth}{0.8pt}}\vspace{2mm}}
\newcommand{\NotesTableStyle}{\linespread{1}\fontsize{9.5pt}{13pt}\selectfont\setlength{\tabcolsep}{4pt}\renewcommand{\arraystretch}{1.1}\setlength{\extrarowheight}{5pt}\setlength{\LTpre}{5pt}\setlength{\LTpost}{6pt}\setlength{\parskip}{0pt}\emergencystretch=2em}
\clubpenalty=10000
\widowpenalty=10000
\displaywidowpenalty=10000

\clearcaptionsetup{table}

% Switch page geometry inside this single document; page numbers stay continuous.
\newcommand{\NotesWanLayout}{%
  \clearpage
  \newgeometry{layoutsize={210mm,297mm},margin=18mm,headheight=15pt}%
  \setlength{\paperwidth}{210mm}%
  \setlength{\paperheight}{297mm}%
  \setlength{\pdfpagewidth}{210mm}%
  \setlength{\pdfpageheight}{297mm}%
  \fancyhfoffset{0pt}%
  \fancyhead[L]{\small\sffamily\color{notesgray}Video VAE / Wan2.1 Details}%
  \setlength{\parskip}{0.26em}%
}

\newcommand{\NotesHThreeLayout}{%
  \clearpage
  \fancyhead[L]{\small\sffamily\color{notesgray}Video VAE / MiniMax-H3 Details}%
}

\newcommand{\NotesLTXLayout}{%
  \clearpage
  \fancyhead[L]{\small\sffamily\color{notesgray}Video VAE / LTX-2.5 Details}%
}

\newcommand{\NotesFluxLayout}{%
  \clearpage
  \fancyhead[L]{\small\sffamily\color{notesgray}Video VAE / FLUX 3 Details}%
  \linespread{1.06}\selectfont
  \setlength{\parskip}{0.24em}%
  \setlength{\abovedisplayskip}{6pt plus 2pt minus 2pt}%
  \setlength{\belowdisplayskip}{6pt plus 2pt minus 2pt}%
  \setlength{\abovedisplayshortskip}{3pt plus 2pt}%
  \setlength{\belowdisplayshortskip}{3pt plus 2pt minus 2pt}%
}

\usepackage{titlesec}
\titleformat{\section}{\Large\bfseries\sffamily\color{notesblue}}{}{0pt}{}
\titleformat{\subsection}{\large\bfseries\sffamily\color{notesblue}}{}{0pt}{}
\titleformat{\subsubsection}{\normalsize\bfseries\sffamily\color{notesblue}}{}{0pt}{}
\titlespacing*{\section}{0pt}{1.5ex}{1ex}
\titlespacing*{\subsection}{0pt}{1.5ex}{1ex}
\titlespacing*{\subsubsection}{0pt}{1.5ex}{0.5ex}
\emergencystretch=2em
\usepackage{bookmark}
\IfFileExists{xurl.sty}{\usepackage{xurl}}{} % add URL line breaks if available
\urlstyle{same}
% fallback for those not using the hyperref driver hyperxmp:
\makeatletter
\@ifundefined{xmpquote}{\newcommand{\xmpquote}[1]{#1}}{}
\makeatother
\hypersetup{
  colorlinks=true,
  linkcolor={Maroon},
  filecolor={Maroon},
  citecolor={Blue},
  urlcolor={Blue},
  pdfcreator={LaTeX via pandoc}}

\author{}
\date{}

\begin{document}

\setlength{\parskip}{3pt}{\sffamily\bfseries\color{notesblue}\fontsize{23}{27}\selectfont VAEs for Video Generation Models\par}\vspace{1mm}{\color{notesblue}\rule{\linewidth}{0.8pt}}\vspace{1mm}

\section{1. Compression Parameters and DiT Token
Counts}\label{compression-parameters-and-dit-token-counts}

Compression factors are listed as \textbf{time × height × width}, and
\(C\) denotes the number of VAE channels. The final column gives the
shape \textbf{after DiT patchification and before projection}, with the
elements within each patch folded into the channel dimension.

\[
R_{\mathrm{nom}}=\frac{3r_t r_h r_w}{C},\qquad
R_{\mathrm{actual}}=\frac{3THW}{C\,T_zH_zW_z},\qquad
N=\frac{T_zH_zW_z}{p_tp_hp_w}.
\]

The DiT patch size is \(p_t\times p_h\times p_w\); patchification only
rearranges elements. Compression ratios are calculated from the VAE
output and already account for patchification inside the VAE. Input
dimensions must be aligned.

\begingroup\NotesTableStyle\fontsize{8.25pt}{9.5pt}\selectfont\renewcommand{\arraystretch}{1.0}\setlength{\extrarowheight}{1pt}

{\def\LTcaptype{none} % do not increment counter
\begin{longtable}[]{@{}
  >{\raggedright\arraybackslash}p{(\linewidth - 12\tabcolsep) * \real{0.1200}}
  >{\raggedright\arraybackslash}p{(\linewidth - 12\tabcolsep) * \real{0.3500}}
  >{\raggedright\arraybackslash}p{(\linewidth - 12\tabcolsep) * \real{0.0850}}
  >{\raggedright\arraybackslash}p{(\linewidth - 12\tabcolsep) * \real{0.0650}}
  >{\raggedright\arraybackslash}p{(\linewidth - 12\tabcolsep) * \real{0.0850}}
  >{\raggedright\arraybackslash}p{(\linewidth - 12\tabcolsep) * \real{0.0850}}
  >{\raggedright\arraybackslash}p{(\linewidth - 12\tabcolsep) * \real{0.2100}}@{}}
\toprule\noalign{}
\begin{minipage}[b]{\linewidth}\raggedright
Model
\end{minipage} & \begin{minipage}[b]{\linewidth}\raggedright
VAE
\end{minipage} & \begin{minipage}[b]{\linewidth}\raggedright
Compression\\
time × height × width\strut
\end{minipage} & \begin{minipage}[b]{\linewidth}\raggedright
Latent\\
channels \(C\)\strut
\end{minipage} & \begin{minipage}[b]{\linewidth}\raggedright
Element compression ratio \(R_{\mathrm{nom}}\)
\end{minipage} & \begin{minipage}[b]{\linewidth}\raggedright
DiT patchify
\end{minipage} & \begin{minipage}[b]{\linewidth}\raggedright
DiT input shape
\end{minipage} \\
\midrule\noalign{}
\endhead
\bottomrule\noalign{}
\endlastfoot
\textbf{MiniMax H3} & \begin{minipage}[t]{\linewidth}\raggedright
\textbf{H3-VisualVAE}\\
Causal 3D CNN encoder + \textbf{non-causal ViT decoder}; reconstructs
spatiotemporal pixel blocks in a single projection.\strut
\end{minipage} & \textbf{\(4\times16\times16\)} & \textbf{24} &
\textbf{128:1} & \(1\times2\times2\) &
\(\left[B,96,T_{\mathrm{H3}},\frac{H}{32},\frac{W}{32}\right]\) \\
\textbf{LTX-2.5} & \begin{minipage}[t]{\linewidth}\raggedright
\textbf{LTX-2.5 Video VAE}\\
Causal 3D CNN encoder + \textbf{non-causal 3D CNN / NA Transformer
diffusion decoder} (CNN in its default configuration); spatial patchify
×4; the diffusion variant uses one-step pixel denoising.\strut
\end{minipage} & \textbf{\(8\times32\times32\)} & \textbf{128} &
\textbf{192:1} & \(1\times1\times1\) &
\(\left[B,128,1+\frac{T-1}{8},\frac{H}{32},\frac{W}{32}\right]\) \\
\textbf{FLUX 3 Action} & \begin{minipage}[t]{\linewidth}\raggedright
\textbf{FLUX 3 Video VAE}\\
Causal NA Transformer encoder + \textbf{non-causal NA Transformer
decoder}; \texttt{5×5×5} neighborhood attention.\strut
\end{minipage} & \textbf{\(4\times32\times32\)} & \textbf{96} &
\textbf{128:1} & \(1\times1\times1\) &
\(\left[B,96,1+\frac{T-1}{4},\frac{H}{32},\frac{W}{32}\right]\) \\
\textbf{Wan 2.1 / 2.2 A14B / Lingbot-Video} &
\begin{minipage}[t]{\linewidth}\raggedright
\textbf{Wan2.1-VAE}\\
Causal 3D CNN encoder + causal 3D CNN decoder; within-frame spatial
attention and per-layer history caches.\strut
\end{minipage} & \textbf{\(4\times8\times8\)} & \textbf{16} &
\textbf{48:1} & \(1\times2\times2\) &
\(\left[B,64,1+\frac{T-1}{4},\frac{H}{16},\frac{W}{16}\right]\) \\
\textbf{Wan 2.2 TI2V-5B} & \begin{minipage}[t]{\linewidth}\raggedright
\textbf{Wan2.2-VAE}\\
Causal 3D CNN encoder + causal 3D CNN decoder; spatial patchify ×2 and
residual shortcuts in resampling modules.\strut
\end{minipage} & \textbf{\(4\times16\times16\)} & \textbf{48} &
\textbf{64:1} & \(1\times2\times2\) &
\(\left[B,192,1+\frac{T-1}{4},\frac{H}{32},\frac{W}{32}\right]\) \\
\textbf{HunyuanVideo-1.5} & \begin{minipage}[t]{\linewidth}\raggedright
\textbf{AutoencoderKLConv3D}\\
Causal 3D CNN encoder + causal 3D CNN decoder; causal attention at the
bottleneck and residual 3D pixel shuffle.\strut
\end{minipage} & \textbf{\(4\times16\times16\)} & \textbf{32} &
\textbf{96:1} & \(1\times1\times1\) &
\(\left[B,32,1+\frac{T-1}{4},\frac{H}{16},\frac{W}{16}\right]\) \\
\textbf{MAGI-2 Preview} & \begin{minipage}[t]{\linewidth}\raggedright
\textbf{Wan2.2-VAE + TurboVAED}\\
Causal 3D CNN encoder + \textbf{non-causal distilled 3D CNN
decoder}.\strut
\end{minipage} & \textbf{\(4\times16\times16\)} & \textbf{48} &
\textbf{64:1} & \(1\times1\times1\) &
\(\left[B,48,1+\frac{T-1}{4},\frac{H}{16},\frac{W}{16}\right]\) \\
\end{longtable}
}

\endgroup

\textbf{H3 handles temporal compression somewhat differently: strictly
speaking, its ratio is not 4:1, and its element compression ratio is
also somewhat lower than 128:1.}

\NotesWanLayout

\section{2. Wan2.1-VAE}\label{wan21}

We now take a closer look at what is probably the most classic VAE:
Wan2.1-VAE. Consider an RGB video with \textbf{81 frames at 480×832}:

\textbf{Full video:} \texttt{{[}3,81,480,832{]}} \(\rightarrow\)
\textbf{latent:} \texttt{{[}16,21,60,104{]}} \(\rightarrow\)
\textbf{reconstructed video:} \texttt{{[}3,81,480,832{]}}.

\[
T_z=1+\frac{81-1}{4}=21,\qquad H_z=\frac{480}{8}=60,\qquad W_z=\frac{832}{8}=104.
\]

\Needspace{8\baselineskip}

\subsection{2.1 Encoder: Tracing a Subsequent 4-Frame Video
Chunk}\label{wan21-encoder}

The full video is divided into \textbf{21 chunks} of \textbf{1, 4, 4,
\ldots{} frames}: one initial single-frame chunk followed by 20
subsequent chunks. Wan2.1-VAE uses a Feature Cache to retain historical
features from earlier chunks. We use one subsequent chunk to trace its
forward pass:

\begingroup\NotesTableStyle\fontsize{10pt}{12pt}\selectfont\renewcommand{\arraystretch}{1.02}\setlength{\extrarowheight}{1pt}

{\def\LTcaptype{none} % do not increment counter
\begin{longtable}[]{@{}
  >{\raggedright\arraybackslash}p{(\linewidth - 4\tabcolsep) * \real{0.1700}}
  >{\raggedright\arraybackslash}p{(\linewidth - 4\tabcolsep) * \real{0.5500}}
  >{\raggedright\arraybackslash}p{(\linewidth - 4\tabcolsep) * \real{0.2800}}@{}}
\toprule\noalign{}
\begin{minipage}[b]{\linewidth}\raggedright
Step
\end{minipage} & \begin{minipage}[b]{\linewidth}\raggedright
Module and operation
\end{minipage} & \begin{minipage}[b]{\linewidth}\raggedright
Output shape
\end{minipage} \\
\midrule\noalign{}
\endhead
\bottomrule\noalign{}
\endlastfoot
Input & Current video chunk & \texttt{{[}3,4,480,832{]}} \\
Entry & Causal convolution, \textbf{3\(\rightarrow\)96} &
\texttt{{[}96,4,480,832{]}} \\
Stage 1 & Residual block ×2, outputting 96 channels \(\rightarrow\)
\textbf{spatial downsampling} & \texttt{{[}96,4,240,416{]}} \\
Stage 2 & Residual block ×2, outputting 192 channels \(\rightarrow\)
\textbf{spatiotemporal downsampling} & \texttt{{[}192,2,120,208{]}} \\
Stage 3 & Residual block ×2, outputting 384 channels \(\rightarrow\)
\textbf{spatiotemporal downsampling} & \texttt{{[}384,1,60,104{]}} \\
Stage 4 & Residual block ×2, outputting 384 channels &
\texttt{{[}384,1,60,104{]}} \\
Bottleneck & Residual block \(\rightarrow\) spatial Attention
\(\rightarrow\) residual block & \texttt{{[}384,1,60,104{]}} \\
Exit & RMSNorm \(\rightarrow\) SiLU \(\rightarrow\) causal convolution,
\textbf{384\(\rightarrow\)32} & \texttt{{[}32,1,60,104{]}} \\
\end{longtable}
}

\endgroup

After all 21 chunks have been encoded, their outputs are concatenated
along the time dimension, passed through a \texttt{1×1×1} convolution,
and split into \(\mu\) and \texttt{log\_var}, each with shape
\texttt{{[}16,21,60,104{]}}. The mean \(\mu\) is standardized per
channel for use as input.

\Needspace{8\baselineskip}

\subsection{2.2 Decoder: Tracing a Subsequent Latent Time
Step}\label{wan21-decoder}

First, \textbf{undo the standardization} of the full latent, then apply
a \textbf{\texttt{1×1×1} convolution, 16\(\rightarrow\)16}; the shape
remains \texttt{{[}16,21,60,104{]}}. Latent temporal positions are then
decoded one at a time. The table traces a subsequent time step:

\begingroup\NotesTableStyle\fontsize{10pt}{12pt}\selectfont\renewcommand{\arraystretch}{1.02}\setlength{\extrarowheight}{1pt}

{\def\LTcaptype{none} % do not increment counter
\begin{longtable}[]{@{}
  >{\raggedright\arraybackslash}p{(\linewidth - 4\tabcolsep) * \real{0.1700}}
  >{\raggedright\arraybackslash}p{(\linewidth - 4\tabcolsep) * \real{0.5500}}
  >{\raggedright\arraybackslash}p{(\linewidth - 4\tabcolsep) * \real{0.2800}}@{}}
\toprule\noalign{}
\begin{minipage}[b]{\linewidth}\raggedright
Step
\end{minipage} & \begin{minipage}[b]{\linewidth}\raggedright
Module and operation
\end{minipage} & \begin{minipage}[b]{\linewidth}\raggedright
Output shape
\end{minipage} \\
\midrule\noalign{}
\endhead
\bottomrule\noalign{}
\endlastfoot
Input & Latent chunk after undoing standardization &
\texttt{{[}16,1,60,104{]}} \\
Entry & Causal convolution, \textbf{16\(\rightarrow\)384} &
\texttt{{[}384,1,60,104{]}} \\
Bottleneck & Residual block \(\rightarrow\) spatial Attention
\(\rightarrow\) residual block & \texttt{{[}384,1,60,104{]}} \\
Stage 1 & Residual block ×3, outputting 384 channels \(\rightarrow\)
\textbf{spatiotemporal upsampling, 384\(\rightarrow\)192} &
\texttt{{[}192,2,120,208{]}} \\
Stage 2 & Residual block ×3, first increasing \textbf{192 to 384}
\(\rightarrow\) \textbf{spatiotemporal upsampling,
384\(\rightarrow\)192} & \texttt{{[}192,4,240,416{]}} \\
Stage 3 & Residual block ×3, outputting 192 channels \(\rightarrow\)
\textbf{spatial upsampling, 192\(\rightarrow\)96} &
\texttt{{[}96,4,480,832{]}} \\
Stage 4 & Residual block ×3, outputting 96 channels &
\texttt{{[}96,4,480,832{]}} \\
Exit & RMSNorm \(\rightarrow\) SiLU \(\rightarrow\) causal convolution,
\textbf{96\(\rightarrow\)3} & \texttt{{[}3,4,480,832{]}} \\
\end{longtable}
}

\endgroup

The first frame skips temporal resampling and produces only \textbf{1
frame}.

\Needspace{8\baselineskip}

\subsection{2.3 Downsampling and Upsampling}\label{wan21-sampling}

\begingroup\NotesTableStyle\fontsize{10pt}{12pt}\selectfont\renewcommand{\arraystretch}{1.02}\setlength{\extrarowheight}{1pt}

{\def\LTcaptype{none} % do not increment counter
\begin{longtable}[]{@{}
  >{\raggedright\arraybackslash}p{(\linewidth - 4\tabcolsep) * \real{0.2000}}
  >{\raggedright\arraybackslash}p{(\linewidth - 4\tabcolsep) * \real{0.5000}}
  >{\raggedright\arraybackslash}p{(\linewidth - 4\tabcolsep) * \real{0.3000}}@{}}
\toprule\noalign{}
\begin{minipage}[b]{\linewidth}\raggedright
Operation
\end{minipage} & \begin{minipage}[b]{\linewidth}\raggedright
Implementation
\end{minipage} & \begin{minipage}[b]{\linewidth}\raggedright
Shape change
\end{minipage} \\
\midrule\noalign{}
\endhead
\bottomrule\noalign{}
\endlastfoot
\textbf{Spatial downsampling} & Pad each frame by 1 on the right and
bottom \(\rightarrow\) \texttt{3×3\ Conv2d}, \texttt{stride=2} & Halve
\(H,W\); \(C,T\) unchanged \\
\textbf{Temporal downsampling} & Concatenate \textbf{1 cached historical
frame} \(\rightarrow\) \texttt{3×1×1} temporal convolution with temporal
stride 2 & Halve \(T\) for subsequent chunks; \(C,H,W\) unchanged \\
\textbf{Spatial upsampling} & Per-frame nearest-neighbor interpolation
×2 (\texttt{nearest-exact}) \(\rightarrow\) \texttt{3×3\ Conv2d},
\(C\rightarrow C/2\) & Double \(H,W\) and halve \(C\); \(T\)
unchanged \\
\textbf{Temporal upsampling} & \texttt{3×1×1} causal convolution,
\(C\rightarrow2C\) \(\rightarrow\) rearrange the two channel groups into
the time dimension & Double \(T\) for subsequent chunks; final \(C,H,W\)
unchanged \\
\end{longtable}
}

\endgroup

\textbf{Execution order:} spatiotemporal downsampling is \textbf{spatial
first, then temporal}; spatiotemporal upsampling is \textbf{temporal
first, then spatial}.

\Needspace{8\baselineskip}

\subsubsection{Residual Blocks, Caches, and
Attention}\label{residual-blocks-caches-and-attention}

\begin{itemize}
\tightlist
\item
  \textbf{Residual blocks:} the main branch applies two sets of
  ``RMSNorm \(\rightarrow\) SiLU \(\rightarrow\) \texttt{3×3×3} causal
  convolution,'' then adds the shortcut. When channel counts differ, the
  shortcut uses a \texttt{1×1×1} convolution to align them. Residual
  blocks preserve \(T,H,W\).
\item
  \textbf{Causal convolution caches:} a standard causal convolution with
  temporal kernel size 3 and stride 1 needs features from the preceding
  \textbf{2 temporal positions}. Missing positions at the beginning of
  the sequence are zero-padded. Each layer maintains its own cache, with
  temporal positions measured at that layer's resolution. Temporal
  downsampling in the table above uses a separate \textbf{1-frame cache}
  path.
\item
  \textbf{Within-frame spatial Attention:} attention operates only on
  the 2D feature map at the same temporal position and preserves
  \texttt{{[}C,T,H,W{]}}. Information across time is conveyed through
  causal convolutions.
\end{itemize}

With Wan2.1-VAE as a reference, we briefly introduce some follow-up
models.

Wan2.2-VAE retains Wan2.1-VAE's causal convolutions and \(4n+1\)
handling, with a temporal compression factor of 4. The main change is
the addition of \texttt{2×2} patchification at the entry, rearranging
spatial positions into the channel dimension. Three subsequent spatial
downsampling stages increase spatial compression from 8×8 to 16×16,
while the latent channel count rises from 16 to 48. Architecturally,
Wan2.2 widens both the encoder and the decoder: the encoder's base
channel count increases from 96 to 160 and its bottleneck from 384 to
640; the decoder's base channel count increases from 96 to 256 and its
bottleneck from 384 to 1024. In addition to the existing residual
blocks, shortcuts are added around entire resampling stages:
downsampling aligns shapes through rearrangement and Avgpool, while
upsampling uses rearrangement and copying, before adding the result to
the main branch.

HunyuanVideo-1.5 keeps 4× temporal compression, increases compression
along both spatial axes from 8× to 16×, and raises the latent channel
count from 16 to 32. Compared with Wan's separate spatial and temporal
resampling, it uses causal convolution with PixelShuffle/Unshuffle-style
rearrangement and adds residual shortcuts to the resampling modules and
the projections on both sides of the latent. Attention in the middle of
the network expands from within-frame spatial attention to
spatiotemporal causal attention, allowing direct attention to current
and historical frames for cross-frame information exchange.

\NotesHThreeLayout

\section{3. MiniMax-H3 VAE}\label{minimax-h3}

Next is MiniMax-H3 VAE. MiniMax H3's default video chunking rule
requires the length to align with \textbf{\(17n+5\)} (H3's handling here
is rather curious; I have considered some possible reasons, but have not
yet found an official explanation). If the original video has \textbf{81
frames (\(4d+1\))}, this example first \textbf{repeats the final frame
to pad the video to 90 frames}, then removes the final 9 frames after
reconstruction.

\textbf{Full video:} \texttt{{[}3,90,480,832{]}} \(\rightarrow\)
\textbf{latent:} \texttt{{[}24,27,30,52{]}} \(\rightarrow\)
\textbf{reconstructed video:} \texttt{{[}3,90,480,832{]}}.

\[
90=17\times5+5,\qquad T_z=5\times5+2=27,\qquad H_z=\frac{480}{16}=30,\qquad W_z=\frac{832}{16}=52.
\]

\Needspace{8\baselineskip}

\subsection{3.1 Encoder: Tracing a 17-Frame Video
Chunk}\label{h3-encoder}

During encoding, the 90-frame video is internally padded again by
repeating the final frame, temporarily reaching \textbf{102 frames}. It
is then split into \textbf{six 17-frame chunks}, each encoded
independently.

\begingroup\NotesTableStyle\fontsize{10pt}{12pt}\selectfont\renewcommand{\arraystretch}{1.02}\setlength{\extrarowheight}{1pt}

{\def\LTcaptype{none} % do not increment counter
\begin{longtable}[]{@{}
  >{\raggedright\arraybackslash}p{(\linewidth - 4\tabcolsep) * \real{0.1700}}
  >{\raggedright\arraybackslash}p{(\linewidth - 4\tabcolsep) * \real{0.5500}}
  >{\raggedright\arraybackslash}p{(\linewidth - 4\tabcolsep) * \real{0.2800}}@{}}
\toprule\noalign{}
\begin{minipage}[b]{\linewidth}\raggedright
Step
\end{minipage} & \begin{minipage}[b]{\linewidth}\raggedright
Module and operation
\end{minipage} & \begin{minipage}[b]{\linewidth}\raggedright
Output shape
\end{minipage} \\
\midrule\noalign{}
\endhead
\bottomrule\noalign{}
\endlastfoot
Input & Current video chunk & \texttt{{[}3,17,480,832{]}} \\
Entry & \texttt{3×3×3} causal convolution, \textbf{3\(\rightarrow\)128}
& \texttt{{[}128,17,480,832{]}} \\
Stage 1 & Residual block ×2, outputting 128 channels \(\rightarrow\)
\textbf{spatial downsampling} & \texttt{{[}128,17,240,416{]}} \\
Stage 2 & Residual block ×2, outputting 256 channels \(\rightarrow\)
\textbf{spatiotemporal downsampling} & \texttt{{[}256,9,120,208{]}} \\
Stage 3 & Residual block ×2, outputting 256 channels \(\rightarrow\)
\textbf{spatiotemporal downsampling} & \texttt{{[}256,5,60,104{]}} \\
Stage 4 & Residual block ×2, outputting 512 channels \(\rightarrow\)
\textbf{spatial downsampling} & \texttt{{[}512,5,30,52{]}} \\
Stage 5 & Residual block ×2, outputting 512 channels &
\texttt{{[}512,5,30,52{]}} \\
Stage 6 & Residual block ×2, outputting 1024 channels &
\texttt{{[}1024,5,30,52{]}} \\
Exit & GroupNorm \(\rightarrow\) SiLU \(\rightarrow\) causal
convolution, \textbf{1024\(\rightarrow\)48} &
\texttt{{[}48,5,30,52{]}} \\
Distribution projection & \texttt{1×1×1} convolution,
\textbf{48\(\rightarrow\)48} & \texttt{{[}48,5,30,52{]}} \\
\end{longtable}
}

\endgroup

Both temporal downsampling steps round up, giving the sequence
\textbf{\(17\rightarrow9\rightarrow5\)}. After all 6 chunks have been
encoded, concatenating along time gives \texttt{{[}48,30,30,52{]}}. With
\texttt{token\_drop=3}, \textbf{3 temporal positions are dropped from
the end of the full sequence}, leaving \texttt{{[}48,27,30,52{]}} (the
discarded portion corresponds to the 12 repeated final frames).

The output is then split into \(\mu\) and \texttt{log\_var}, each with
shape \texttt{{[}24,27,30,52{]}}. The latent is sampled from the
posterior distribution and then standardized per channel.

\Needspace{8\baselineskip}

\subsection{3.2 Decoder: Tracing a Latent Window}\label{h3-decoder}

First, \textbf{undo the standardization} of the full latent. Decode it
in \textbf{windows of 7 temporal positions, advancing by 5 positions
each time}, so neighboring windows overlap by 2 positions. We trace one
window of shape \texttt{{[}24,7,30,52{]}}; its video token count and
hidden dimension are:

\[
N=7\times30\times52=10920,\qquad D=32\times64=2048.
\]

\begingroup\NotesTableStyle\fontsize{10pt}{12pt}\selectfont\renewcommand{\arraystretch}{1.02}\setlength{\extrarowheight}{1pt}

{\def\LTcaptype{none} % do not increment counter
\begin{longtable}[]{@{}
  >{\raggedright\arraybackslash}p{(\linewidth - 4\tabcolsep) * \real{0.1700}}
  >{\raggedright\arraybackslash}p{(\linewidth - 4\tabcolsep) * \real{0.5500}}
  >{\raggedright\arraybackslash}p{(\linewidth - 4\tabcolsep) * \real{0.2800}}@{}}
\toprule\noalign{}
\begin{minipage}[b]{\linewidth}\raggedright
Step
\end{minipage} & \begin{minipage}[b]{\linewidth}\raggedright
Module and operation
\end{minipage} & \begin{minipage}[b]{\linewidth}\raggedright
Output shape
\end{minipage} \\
\midrule\noalign{}
\endhead
\bottomrule\noalign{}
\endlastfoot
Input & Latent window after undoing standardization &
\texttt{{[}24,7,30,52{]}} \\
Entry & \texttt{1×1×1} convolution, \textbf{24\(\rightarrow\)24} &
\texttt{{[}24,7,30,52{]}} \\
Flatten & Turn each spatiotemporal position into one token &
\texttt{{[}10920,24{]}} \\
Feature projection & Linear, \textbf{24\(\rightarrow\)2048} &
\texttt{{[}10920,2048{]}} \\
Extra tokens & Append \textbf{4 register tokens + 1 initially zero
token} (a classic CV move) & \texttt{{[}10925,2048{]}} \\
Backbone & Transformer block ×36, 3D RoPE & \texttt{{[}10925,2048{]}} \\
Exit & LayerNorm \(\rightarrow\) Linear,
\textbf{2048\(\rightarrow\)3072} & \texttt{{[}10925,3072{]}} \\
Remove extra tokens & Keep the video tokens &
\texttt{{[}10920,3072{]}} \\
Pixel rearrangement & Expand each token into a \textbf{\texttt{4×16×16}
RGB block} & \texttt{{[}3,28,480,832{]}} \\
\end{longtable}
}

\endgroup

The output projection maps each video token to
\textbf{\(3072=3\times4\times16\times16\)} values, then rearranges them
into a \texttt{4×16×16} RGB block. A window therefore first produces
\texttt{{[}3,28,480,832{]}}, followed by temporal cropping and window
blending.

\Needspace{8\baselineskip}

\subsubsection{Temporal Cropping and Full-Sequence
Stitching}\label{h3-time-stitch}

Each window's 28 output frames are split into the first \textbf{20
frames} and the final \textbf{8 frames}. The first 3 frames of each part
are removed, leaving a \textbf{17-frame main segment + a 5-frame overlap
region}.

The 27 latent temporal positions form \textbf{5 windows}. With
zero-based indexing, their ranges are \texttt{0-6}, \texttt{5-11},
\texttt{10-16}, \texttt{15-21}, and \texttt{20-26}. Each window retains
22 frames, and adjacent windows blend their 5 overlapping frames,
yielding:

\[
5\times22-4\times5=90\text{ frames}.
\]

The blended output has shape \texttt{{[}3,90,480,832{]}}. Removing the
final 9 frames added in this example restores
\texttt{{[}3,81,480,832{]}}.

\Needspace{8\baselineskip}

\subsection{3.3 Downsampling and Pixel Expansion}\label{h3-operations}

\begingroup\NotesTableStyle\fontsize{10pt}{12pt}\selectfont\renewcommand{\arraystretch}{1.02}\setlength{\extrarowheight}{1pt}

{\def\LTcaptype{none} % do not increment counter
\begin{longtable}[]{@{}
  >{\raggedright\arraybackslash}p{(\linewidth - 4\tabcolsep) * \real{0.2000}}
  >{\raggedright\arraybackslash}p{(\linewidth - 4\tabcolsep) * \real{0.5000}}
  >{\raggedright\arraybackslash}p{(\linewidth - 4\tabcolsep) * \real{0.3000}}@{}}
\toprule\noalign{}
\begin{minipage}[b]{\linewidth}\raggedright
Operation
\end{minipage} & \begin{minipage}[b]{\linewidth}\raggedright
Implementation
\end{minipage} & \begin{minipage}[b]{\linewidth}\raggedright
Shape change
\end{minipage} \\
\midrule\noalign{}
\endhead
\bottomrule\noalign{}
\endlastfoot
\textbf{Spatial downsampling} & Reflection-pad by 1 on the right and
bottom \(\rightarrow\) \texttt{3×3×3} causal convolution, stride
\texttt{(1,2,2)} & Halve \(H,W\); \(C,T\) unchanged \\
\textbf{Spatiotemporal downsampling} & Reflection-pad by 1 on the right
and bottom \(\rightarrow\) \texttt{3×3×3} causal convolution, stride
\texttt{(2,2,2)} & Halve \(H,W\); \(T\to\lceil T/2\rceil\); \(C\)
unchanged \\
\textbf{Pixel expansion} & Linear \(\rightarrow\) spatiotemporal
rearrangement & Each video token produces one \texttt{4×16×16} RGB
block \\
\end{longtable}
}

\endgroup

\textbf{Downsampling order:} spatial \(\rightarrow\) spatiotemporal
\(\rightarrow\) spatiotemporal \(\rightarrow\) spatial, at the ends of
the first four stages, respectively. Temporal convolutions in both
downsampling types use 2 zero-padded positions on the left.

\Needspace{8\baselineskip}

\subsubsection{Residual Blocks, Normalization, and
Attention}\label{residual-blocks-normalization-and-attention}

\begin{itemize}
\tightlist
\item
  \textbf{Residual blocks:} the main branch applies two sets of
  ``GroupNorm \(\rightarrow\) SiLU \(\rightarrow\) \texttt{3×3×3} causal
  convolution,'' then adds the shortcut. When channel counts differ, the
  shortcut uses a \texttt{1×1×1} convolution to align them. Residual
  blocks preserve \(T,H,W\).
\item
  \textbf{Causal encoding:} temporal convolutions use zero-padding only
  on the left, and GroupNorm is computed independently at each temporal
  position. Each 17-frame chunk enters the Encoder independently.
\item
  \textbf{ViT decoder blocks:} each block applies RMSNorm
  \(\rightarrow\) Attention, followed by RMSNorm \(\rightarrow\) gated
  SiLU FFN. Both branches are scaled by learnable parameters before
  being added back to the residual, preserving \texttt{{[}N,2048{]}}.
\item
  \textbf{Attention scope:} the ViT Decoder in the released
  configuration is non-causal and can attend to all temporal positions
  within the current window.
\end{itemize}

\NotesLTXLayout

\section{4. LTX-2.5 Video VAE}\label{ltx25}

Next is LTX-2.5's video VAE. It uses a \textbf{causal 3D CNN Encoder}
and offers two Decoders: a \textbf{non-causal 3D CNN decoder} and a
\textbf{non-causal NA Transformer diffusion decoder}. Both share the
same latent space, with a temporal compression factor of \textbf{8},
spatial compression factors of \textbf{32} along both height and width,
and \textbf{128} latent channels.

Again, consider an RGB video with \textbf{81 frames at 480×832}. Shapes
are written as \texttt{{[}C,T,H,W{]}}, omitting \texttt{B=1}:

\textbf{Full video:} \texttt{{[}3,81,480,832{]}} \(\rightarrow\)
\textbf{latent:} \texttt{{[}128,11,15,26{]}} \(\rightarrow\)
\textbf{reconstructed video:} \texttt{{[}3,81,480,832{]}}.

\[
T_z=1+\frac{81-1}{8}=11,\qquad H_z=\frac{480}{32}=15,\qquad W_z=\frac{832}{32}=26.
\]

We trace the forward pass of the full video below. Chunking and overlap
blending are separate execution strategies; we do not treat the video as
11 independently encoded and decoded chunks from the outset.

\Needspace{8\baselineskip}

\subsection{4.1 Encoder: Tracing the Full 81-Frame
Video}\label{ltx25-encoder}

The entry first applies \textbf{4×4 spatial patchification}, rearranging
the 16 neighboring pixel positions in each frame into the channel
dimension while keeping the temporal length unchanged. Four downsampling
stages follow: \textbf{spatial \(\rightarrow\) temporal \(\rightarrow\)
spatiotemporal \(\rightarrow\) spatiotemporal}.

\begingroup\NotesTableStyle\fontsize{10pt}{12pt}\selectfont\renewcommand{\arraystretch}{1.02}\setlength{\extrarowheight}{1pt}

{\def\LTcaptype{none} % do not increment counter
\begin{longtable}[]{@{}
  >{\raggedright\arraybackslash}p{(\linewidth - 4\tabcolsep) * \real{0.1700}}
  >{\raggedright\arraybackslash}p{(\linewidth - 4\tabcolsep) * \real{0.5500}}
  >{\raggedright\arraybackslash}p{(\linewidth - 4\tabcolsep) * \real{0.2800}}@{}}
\toprule\noalign{}
\begin{minipage}[b]{\linewidth}\raggedright
Step
\end{minipage} & \begin{minipage}[b]{\linewidth}\raggedright
Module and operation
\end{minipage} & \begin{minipage}[b]{\linewidth}\raggedright
Output shape
\end{minipage} \\
\midrule\noalign{}
\endhead
\bottomrule\noalign{}
\endlastfoot
Input & RGB video & \texttt{{[}3,81,480,832{]}} \\
Patchify & Spatial \texttt{4×4} rearrangement,
\textbf{3\(\rightarrow\)48} & \texttt{{[}48,81,120,208{]}} \\
Entry & \texttt{3×3×3} causal convolution, \textbf{48\(\rightarrow\)128}
& \texttt{{[}128,81,120,208{]}} \\
Stage 1 & Residual block ×4, 128 channels \(\rightarrow\)
\textbf{spatial downsampling, 128\(\rightarrow\)256} &
\texttt{{[}256,81,60,104{]}} \\
Stage 2 & Residual block ×6, 256 channels \(\rightarrow\)
\textbf{temporal downsampling, 256\(\rightarrow\)512} &
\texttt{{[}512,41,60,104{]}} \\
Stage 3 & Residual block ×4, 512 channels \(\rightarrow\)
\textbf{spatiotemporal downsampling, 512\(\rightarrow\)1024} &
\texttt{{[}1024,21,30,52{]}} \\
Stage 4 & Residual block ×2, 1024 channels \(\rightarrow\)
\textbf{spatiotemporal downsampling, 1024\(\rightarrow\)1024} &
\texttt{{[}1024,11,15,26{]}} \\
Bottleneck & Residual block ×2, 1024 channels &
\texttt{{[}1024,11,15,26{]}} \\
Exit & PixelNorm \(\rightarrow\) SiLU \(\rightarrow\) causal
convolution, \textbf{1024\(\rightarrow\)129} &
\texttt{{[}129,11,15,26{]}} \\
\end{longtable}
}

\endgroup

All three temporal downsampling steps preserve the position
corresponding to the first frame. The temporal lengths are:

\[
81\rightarrow41\rightarrow21\rightarrow11.
\]

The first \textbf{128 channels of the convolutional head's output are
\(\mu\)}. As in the other examples, normalization using per-channel
means and standard deviations gives \texttt{{[}128,11,15,26{]}}. Unlike
the other VAEs, this one has no additional \texttt{1×1×1} distribution
projection convolution.

\clearpage

\subsection{4.2 CNN Decoder: Reconstructing the Video Stage by
Stage}\label{ltx25-conv-decoder}

First undo the latent's per-channel normalization, then apply
convolutional residual blocks and four upsampling stages. Finally,
unpatchify restores the RGB pixels.

\begingroup\NotesTableStyle\fontsize{10pt}{12pt}\selectfont\renewcommand{\arraystretch}{1.02}\setlength{\extrarowheight}{1pt}

{\def\LTcaptype{none} % do not increment counter
\begin{longtable}[]{@{}
  >{\raggedright\arraybackslash}p{(\linewidth - 4\tabcolsep) * \real{0.1700}}
  >{\raggedright\arraybackslash}p{(\linewidth - 4\tabcolsep) * \real{0.5500}}
  >{\raggedright\arraybackslash}p{(\linewidth - 4\tabcolsep) * \real{0.2800}}@{}}
\toprule\noalign{}
\begin{minipage}[b]{\linewidth}\raggedright
Step
\end{minipage} & \begin{minipage}[b]{\linewidth}\raggedright
Module and operation
\end{minipage} & \begin{minipage}[b]{\linewidth}\raggedright
Output shape
\end{minipage} \\
\midrule\noalign{}
\endhead
\bottomrule\noalign{}
\endlastfoot
Input & Latent after undoing normalization &
\texttt{{[}128,11,15,26{]}} \\
Entry & \texttt{3×3×3} convolution, \textbf{128\(\rightarrow\)1024} &
\texttt{{[}1024,11,15,26{]}} \\
Stage 1 & Residual block ×2, 1024 channels \(\rightarrow\)
\textbf{spatiotemporal upsampling, 1024\(\rightarrow\)512} &
\texttt{{[}512,21,30,52{]}} \\
Stage 2 & Residual block ×2, 512 channels \(\rightarrow\)
\textbf{spatiotemporal upsampling, 512\(\rightarrow\)512} &
\texttt{{[}512,41,60,104{]}} \\
Stage 3 & Residual block ×4, 512 channels \(\rightarrow\)
\textbf{temporal upsampling, 512\(\rightarrow\)256} &
\texttt{{[}256,81,60,104{]}} \\
Stage 4 & Residual block ×6, 256 channels \(\rightarrow\)
\textbf{spatial upsampling, 256\(\rightarrow\)128} &
\texttt{{[}128,81,120,208{]}} \\
Stage 5 & Residual block ×4, 128 channels &
\texttt{{[}128,81,120,208{]}} \\
Exit & PixelNorm \(\rightarrow\) SiLU \(\rightarrow\) \texttt{3×3×3}
convolution, \textbf{128\(\rightarrow\)48} &
\texttt{{[}48,81,120,208{]}} \\
Unpatchify & Rearrange channels into \texttt{4×4} spatial pixel blocks,
\textbf{48\(\rightarrow\)3} & \texttt{{[}3,81,480,832{]}} \\
\end{longtable}
}

\endgroup

Temporal upsampling first doubles the length and then removes the first
temporal position, so each step follows \textbf{\(T\rightarrow2T-1\)}:

\[
11\rightarrow21\rightarrow41\rightarrow81.
\]

The convolutions in this Decoder run in \textbf{non-causal mode},
allowing them to use information from both preceding and subsequent
temporal positions.

\Needspace{8\baselineskip}

\subsection{4.3 Diffusion Decoder: Conditioning Features and Pixel
Denoising}\label{ltx25-diffusion-decoder}

This decoder first upsamples the latent into conditioning features, or
\textbf{context}. Conditioned on this context, it reconstructs the video
from pixel noise in \textbf{1 denoising step, directly predicting the
reconstructed video \(x_0\)}.

\Needspace{8\baselineskip}

\subsubsection{Generating Conditioning Features}\label{ltx25-context}

The official implementation repeats the final latent temporal position
\textbf{2 additional times} to handle the end boundary of neighborhood
Attention. We include this temporary padding in the shapes below and
crop it away at the end.

\begingroup\NotesTableStyle\fontsize{10pt}{12pt}\selectfont\renewcommand{\arraystretch}{1.02}\setlength{\extrarowheight}{1pt}

{\def\LTcaptype{none} % do not increment counter
\begin{longtable}[]{@{}
  >{\raggedright\arraybackslash}p{(\linewidth - 4\tabcolsep) * \real{0.1700}}
  >{\raggedright\arraybackslash}p{(\linewidth - 4\tabcolsep) * \real{0.5500}}
  >{\raggedright\arraybackslash}p{(\linewidth - 4\tabcolsep) * \real{0.2800}}@{}}
\toprule\noalign{}
\begin{minipage}[b]{\linewidth}\raggedright
Step
\end{minipage} & \begin{minipage}[b]{\linewidth}\raggedright
Module and operation
\end{minipage} & \begin{minipage}[b]{\linewidth}\raggedright
Output shape
\end{minipage} \\
\midrule\noalign{}
\endhead
\bottomrule\noalign{}
\endlastfoot
Input & Latent after undoing normalization &
\texttt{{[}128,11,15,26{]}} \\
Boundary handling & Repeat the final latent temporal position 2
additional times & \texttt{{[}128,13,15,26{]}} \\
Entry & Position-wise Linear, \textbf{128\(\rightarrow\)2048} &
\texttt{{[}2048,13,15,26{]}} \\
Stage 1 & NA Block ×4 \(\rightarrow\) \textbf{spatial upsampling,
2048\(\rightarrow\)1024} & \texttt{{[}1024,13,30,52{]}} \\
Stage 2 & NA Block ×6 \(\rightarrow\) \textbf{temporal upsampling,
1024\(\rightarrow\)512} & \texttt{{[}512,25,30,52{]}} \\
Stage 3 & NA Block ×4 \(\rightarrow\) \textbf{spatiotemporal upsampling,
512\(\rightarrow\)512} & \texttt{{[}512,49,60,104{]}} \\
Stage 4 & NA Block ×2 \(\rightarrow\) \textbf{spatiotemporal upsampling,
512\(\rightarrow\)256} & \texttt{{[}256,97,120,208{]}} \\
Crop & Remove the final \textbf{16 temporal positions} corresponding to
the temporary padding & \texttt{{[}256,81,120,208{]}} \\
\end{longtable}
}

\endgroup

The 2 added latent temporal positions undergo 8× temporal expansion,
corresponding to \textbf{16 temporal positions}, so cropping gives
\textbf{\(97-16=81\)}. The resulting context already has the target
video's temporal length, while its spatial resolution remains
\textbf{1/4×1/4} of the target.

\clearpage

\subsubsection{Pixel Denoising and
Reconstruction}\label{ltx25-pixel-decode}

Initialize random pixel noise with the same shape as the target video.
We trace the pixel feature branch below; context is injected as
conditioning in each diffusion block.

\begingroup\NotesTableStyle\fontsize{10pt}{12pt}\selectfont\renewcommand{\arraystretch}{1.02}\setlength{\extrarowheight}{1pt}

{\def\LTcaptype{none} % do not increment counter
\begin{longtable}[]{@{}
  >{\raggedright\arraybackslash}p{(\linewidth - 4\tabcolsep) * \real{0.1700}}
  >{\raggedright\arraybackslash}p{(\linewidth - 4\tabcolsep) * \real{0.5500}}
  >{\raggedright\arraybackslash}p{(\linewidth - 4\tabcolsep) * \real{0.2800}}@{}}
\toprule\noalign{}
\begin{minipage}[b]{\linewidth}\raggedright
Step
\end{minipage} & \begin{minipage}[b]{\linewidth}\raggedright
Module and operation
\end{minipage} & \begin{minipage}[b]{\linewidth}\raggedright
Output shape
\end{minipage} \\
\midrule\noalign{}
\endhead
\bottomrule\noalign{}
\endlastfoot
Input & Random pixel noise \(x_t\) & \texttt{{[}3,81,480,832{]}} \\
Patchify & Spatial \texttt{4×4} rearrangement,
\textbf{3\(\rightarrow\)48} & \texttt{{[}48,81,120,208{]}} \\
Entry & Position-wise Linear, \textbf{48\(\rightarrow\)256} &
\texttt{{[}256,81,120,208{]}} \\
Denoising backbone & Diffusion NA Block ×8, injecting context and
denoising timestep conditioning & \texttt{{[}256,81,120,208{]}} \\
Exit & RMSNorm \(\rightarrow\) Linear, \textbf{256\(\rightarrow\)48} &
\texttt{{[}48,81,120,208{]}} \\
Unpatchify & Rearrange channels into \texttt{4×4} RGB pixel blocks,
\textbf{48\(\rightarrow\)3} & \texttt{{[}3,81,480,832{]}} \\
\end{longtable}
}

\endgroup

The final output is the reconstructed video \(x_0\). The \textbf{8
blocks refer to network depth}; the ``denoising timestep'' indicates the
noise level and is distinct from the video's 81 temporal positions.

\Needspace{8\baselineskip}

\subsection{4.4 NA Attention: Local Spatiotemporal
Attention}\label{ltx25-na}

\textbf{NA stands for Neighborhood Attention.} At each feature-map
position \texttt{(t,h,w)}, it takes only the K and V values within a
nearby 3D window and computes Attention with the Q at the current
position. The window slides with the query position, and neighboring
windows overlap.

For example, away from boundaries, a \textbf{\texttt{3×7×7}} window
covers the \textbf{preceding, current, and following temporal
positions}, taking \textbf{7×7} spatial grid points at each one. Each
query in each attention head therefore attends to
\textbf{\(3\times7\times7=147\)} positions, including itself.

\begingroup\NotesTableStyle\fontsize{10pt}{12pt}\selectfont\renewcommand{\arraystretch}{1.02}\setlength{\extrarowheight}{1pt}

{\def\LTcaptype{none} % do not increment counter
\begin{longtable}[]{@{}
  >{\raggedright\arraybackslash}p{(\linewidth - 10\tabcolsep) * \real{0.2200}}
  >{\raggedright\arraybackslash}p{(\linewidth - 10\tabcolsep) * \real{0.1200}}
  >{\raggedright\arraybackslash}p{(\linewidth - 10\tabcolsep) * \real{0.1500}}
  >{\raggedright\arraybackslash}p{(\linewidth - 10\tabcolsep) * \real{0.1400}}
  >{\raggedright\arraybackslash}p{(\linewidth - 10\tabcolsep) * \real{0.2300}}
  >{\raggedright\arraybackslash}p{(\linewidth - 10\tabcolsep) * \real{0.1400}}@{}}
\toprule\noalign{}
\begin{minipage}[b]{\linewidth}\raggedright
Stage
\end{minipage} & \begin{minipage}[b]{\linewidth}\raggedright
Feature channels
\end{minipage} & \begin{minipage}[b]{\linewidth}\raggedright
Attention heads
\end{minipage} & \begin{minipage}[b]{\linewidth}\raggedright
Head dimension
\end{minipage} & \begin{minipage}[b]{\linewidth}\raggedright
Window: time × height × width
\end{minipage} & \begin{minipage}[b]{\linewidth}\raggedright
Positions per head
\end{minipage} \\
\midrule\noalign{}
\endhead
\bottomrule\noalign{}
\endlastfoot
Stage 1 & 2048 & 32 & 64 & \texttt{3×7×7} & 147 \\
Stage 2 & 1024 & 16 & 64 & \texttt{3×7×7} & 147 \\
Stage 3 & 512 & 8 & 64 & \texttt{3×5×5} & 75 \\
Stage 4 & 512 & 8 & 64 & \texttt{3×5×5} & 75 \\
Stage 5: pixel denoising & 256 & 4 & 64 & \texttt{11×11×11} & 1331 \\
\end{longtable}
}

\endgroup

Every stage uses a head dimension of \textbf{64}, with \textbf{\(C/64\)}
heads. Window sizes are measured in \textbf{feature grid points at the
corresponding layer}. NA itself preserves \texttt{{[}C,T,H,W{]}};
separate modules perform upsampling.

A standard NA Block applies:

\begin{verbatim}
RMSNorm -> 3D NA Attention -> add residual
RMSNorm -> SwiGLU MLP      -> add residual
\end{verbatim}

Q and K use per-head RMSNorm and 3D RoPE; the MLP hidden dimension is
\textbf{4C}. The Diffusion NA Blocks in Stage 5 additionally inject
context and modulate features using the denoising timestep.

All these Attention layers are \textbf{non-causal}, allowing attention
to subsequent temporal positions. At boundaries, the window shifts into
the valid region to retain its size. For example, with a temporal window
of 3, the first position attends to \texttt{0,\ 1,\ 2}.

\clearpage

\subsection{4.5 Downsampling, Upsampling, and Temporal
Boundaries}\label{ltx25-sampling}

\begingroup\NotesTableStyle\fontsize{10pt}{12pt}\selectfont\renewcommand{\arraystretch}{1.02}\setlength{\extrarowheight}{1pt}

{\def\LTcaptype{none} % do not increment counter
\begin{longtable}[]{@{}
  >{\raggedright\arraybackslash}p{(\linewidth - 4\tabcolsep) * \real{0.2000}}
  >{\raggedright\arraybackslash}p{(\linewidth - 4\tabcolsep) * \real{0.5000}}
  >{\raggedright\arraybackslash}p{(\linewidth - 4\tabcolsep) * \real{0.3000}}@{}}
\toprule\noalign{}
\begin{minipage}[b]{\linewidth}\raggedright
Operation
\end{minipage} & \begin{minipage}[b]{\linewidth}\raggedright
Implementation
\end{minipage} & \begin{minipage}[b]{\linewidth}\raggedright
Shape change
\end{minipage} \\
\midrule\noalign{}
\endhead
\bottomrule\noalign{}
\endlastfoot
\textbf{Spatial Patchify} & Rearrange each frame's \texttt{4×4} pixel
blocks into the channel dimension & Divide \(H,W\) by 4; multiply \(C\)
by 16; \(T\) unchanged \\
\textbf{Spatial downsampling} & Causal convolution \(\rightarrow\)
rearrange spatial positions into channels; add a shortcut obtained by
rearrangement and grouped averaging & Halve \(H,W\); \(T\) unchanged \\
\textbf{Temporal downsampling} & Repeat the first frame once
\(\rightarrow\) causal convolution and temporal rearrangement; add a
shortcut & \(T\rightarrow(T+1)/2\); \(H,W\) unchanged \\
\textbf{Spatiotemporal downsampling} & Repeat the first frame once
\(\rightarrow\) causal convolution and joint \texttt{2×2×2}
rearrangement; add a shortcut & \(T\rightarrow(T+1)/2\); halve
\(H,W\) \\
\textbf{CNN upsampling} & Convolution produces the required channels
\(\rightarrow\) rearrange channels into the specified spatial/temporal
dimensions & Double the specified axes; remove the first frame after
temporal expansion \\
\textbf{NA Decoder upsampling} & Linear produces the required channels
\(\rightarrow\) rearrange channels into the specified spatial/temporal
dimensions & Double the specified axes; remove the first frame after
temporal expansion \\
\textbf{Spatial Unpatchify} & Rearrange channels back into each frame's
\texttt{4×4} pixel blocks & Multiply \(H,W\) by 4; divide \(C\) by 16;
\(T\) unchanged \\
\end{longtable}
}

\endgroup

The main downsampling branch first adjusts channels through convolution,
then folds local spatiotemporal positions into the channel dimension.
The shortcut applies the same rearrangement to the input and uses
grouped averaging to match the output channel count. Spatiotemporal
resampling uses joint rearrangement.

\Needspace{8\baselineskip}

\subsubsection{Residual Blocks, Normalization, and
Chunking}\label{residual-blocks-normalization-and-chunking}

\begin{itemize}
\tightlist
\item
  \textbf{CNN residual blocks:} the main branch primarily applies two
  sets of ``PixelNorm \(\rightarrow\) SiLU \(\rightarrow\)
  \texttt{3×3×3} convolution,'' then adds the shortcut, preserving
  \(T,H,W\).
\item
  \textbf{PixelNorm:} RMS normalization along the channel dimension at
  each spatiotemporal position.
\item
  \textbf{Temporal boundaries:} the Encoder's causal convolutions use
  first-frame replication for left temporal padding. Rearrangement
  modules that include temporal downsampling also repeat the first frame
  once more. The Decoder removes leading frames after temporal
  expansion, giving the overall length
  \textbf{\(T_{\mathrm{out}}=8(T_z-1)+1\)}.
\item
  \textbf{Chunked execution:} LTX supports tiling with overlap regions.
  Here, \textbf{\(81=1+10\times8\)} expresses the temporal compression
  relationship. Actual decoding must retain context across time; the 11
  latent temporal positions cannot simply be decoded independently and
  concatenated.
\end{itemize}

\NotesFluxLayout

\section{5. FLUX 3 Video VAE}\label{flux3}

Next is the video VAE in FLUX 3 Action. It uses a \textbf{causal NA
Transformer Encoder + non-causal NA Transformer Decoder}, with a
temporal compression factor of \textbf{4}, spatial compression factors
of \textbf{32} along both height and width, and \textbf{96} latent
channels.

Again, consider an RGB video with \textbf{81 frames at 480×832}. Shapes
are written as \texttt{{[}C,T,H,W{]}}, omitting \texttt{B=1}:

\textbf{Full video:} \texttt{{[}3,81,480,832{]}} \(\rightarrow\)
\textbf{latent:} \texttt{{[}96,21,15,26{]}} \(\rightarrow\)
\textbf{reconstructed video:} \texttt{{[}3,81,480,832{]}}.

\[
T_z=1+\frac{81-1}{4}=21,\qquad H_z=\frac{480}{32}=15,\qquad W_z=\frac{832}{32}=26.
\]

\Needspace{6\baselineskip}

\subsection{5.1 Blocks and Attention}\label{flux3-attention}

The FLUX 3 VAE also uses Neighborhood Attention, with
\textbf{\texttt{5×5×5}} windows throughout, together with QK Norm and 3D
RoPE. \textbf{The Encoder uses temporally causal Attention, while the
Decoder uses non-causal Attention.}

\begingroup\NotesTableStyle\fontsize{10pt}{12pt}\selectfont\renewcommand{\arraystretch}{1.02}\setlength{\extrarowheight}{1pt}

{\def\LTcaptype{none} % do not increment counter
\begin{longtable}[]{@{}
  >{\raggedright\arraybackslash}p{(\linewidth - 6\tabcolsep) * \real{0.2000}}
  >{\raggedright\arraybackslash}p{(\linewidth - 6\tabcolsep) * \real{0.2500}}
  >{\raggedright\arraybackslash}p{(\linewidth - 6\tabcolsep) * \real{0.2000}}
  >{\raggedright\arraybackslash}p{(\linewidth - 6\tabcolsep) * \real{0.3500}}@{}}
\toprule\noalign{}
\begin{minipage}[b]{\linewidth}\raggedright
Feature channels
\end{minipage} & \begin{minipage}[b]{\linewidth}\raggedright
Attention heads
\end{minipage} & \begin{minipage}[b]{\linewidth}\raggedright
Head dimension
\end{minipage} & \begin{minipage}[b]{\linewidth}\raggedright
Window: time × height × width
\end{minipage} \\
\midrule\noalign{}
\endhead
\bottomrule\noalign{}
\endlastfoot
256 & 4 & 64 & \texttt{5×5×5} \\
512 & 8 & 64 & \texttt{5×5×5} \\
1024 & 16 & 64 & \texttt{5×5×5} \\
2048 & 32 & 64 & \texttt{5×5×5} \\
\end{longtable}
}

\endgroup

Within a single NA layer, the Encoder attends only to \textbf{the
current temporal position and up to 4 preceding positions}. Away from
boundaries, the Decoder attends to \textbf{the preceding 2 positions,
the current position, and the following 2 positions}; at boundaries, the
window shifts toward the sequence interior. Single-frame inputs use
\textbf{\texttt{5×5} 2D NA}.

\Needspace{6\baselineskip}

\subsection{5.2 Actual Chunking and Context}\label{flux3-chunking}

\textbf{Encoding chunks:} execution uses \textbf{45-frame chunks,
advancing by 44 frames each time}. Neighboring chunks overlap by 1
frame, and each chunk enters the Encoder independently. For the classic
81-frame example, the final frame is first repeated to \textbf{pad 81
frames to 89}:

\begin{verbatim}
Chunk 1: frames 1-45
Chunk 2: frames 45-89
\end{verbatim}

Each chunk produces \textbf{12 latent temporal positions}. Keep the
entire output of the first chunk, discard the first latent temporal
position from subsequent chunks, then concatenate and crop according to
the original video length:

\[
12+(12-1)=23
\quad\longrightarrow\quad
\text{keep the first }21\text{ temporal positions}.
\]

\Needspace{6\baselineskip}

\subsection{5.3 Encoder: Tracing a 45-Frame Video
Chunk}\label{flux3-encoder}

The entry divides pixels into \textbf{\texttt{1×4×4}} patches and
projects them to 256 channels, implemented as a Conv3d with both kernel
and stride set to \texttt{1×4×4}. Three spatial merging steps follow,
then two temporal merging steps at low spatial resolution.

\begingroup\NotesTableStyle\fontsize{10pt}{12pt}\selectfont\renewcommand{\arraystretch}{1.02}\setlength{\extrarowheight}{1pt}

{\def\LTcaptype{none} % do not increment counter
\begin{longtable}[]{@{}
  >{\raggedright\arraybackslash}p{(\linewidth - 4\tabcolsep) * \real{0.1700}}
  >{\raggedright\arraybackslash}p{(\linewidth - 4\tabcolsep) * \real{0.5500}}
  >{\raggedright\arraybackslash}p{(\linewidth - 4\tabcolsep) * \real{0.2800}}@{}}
\toprule\noalign{}
\begin{minipage}[b]{\linewidth}\raggedright
Step
\end{minipage} & \begin{minipage}[b]{\linewidth}\raggedright
Module and operation
\end{minipage} & \begin{minipage}[b]{\linewidth}\raggedright
Output shape
\end{minipage} \\
\midrule\noalign{}
\endhead
\bottomrule\noalign{}
\endlastfoot
Input & Current video chunk & \texttt{{[}3,45,480,832{]}} \\
Entry & \texttt{1×4×4\ Conv3d}, stride \texttt{(1,4,4)},
\textbf{3\(\rightarrow\)256} & \texttt{{[}256,45,120,208{]}} \\
Stage 1 & NA Block ×1, 256 channels \(\rightarrow\) \textbf{spatial
merging, 256\(\rightarrow\)512} & \texttt{{[}512,45,60,104{]}} \\
Stage 2 & NA Block ×4, 512 channels \(\rightarrow\) \textbf{spatial
merging, 512\(\rightarrow\)1024} & \texttt{{[}1024,45,30,52{]}} \\
Stage 3 backbone & NA Block ×8, 1024 channels \(\rightarrow\)
\textbf{spatial merging, 1024\(\rightarrow\)2048} &
\texttt{{[}2048,45,15,26{]}} \\
First temporal compression & Additional NA Block ×1 \(\rightarrow\)
\textbf{temporal merging, 45\(\rightarrow\)23} &
\texttt{{[}2048,23,15,26{]}} \\
Stage 4 backbone & NA Block ×8, 2048 channels &
\texttt{{[}2048,23,15,26{]}} \\
Second temporal compression & Additional NA Block ×1 \(\rightarrow\)
\textbf{temporal merging, 23\(\rightarrow\)12} &
\texttt{{[}2048,12,15,26{]}} \\
Exit & Position-wise Linear, \textbf{2048\(\rightarrow\)192} &
\texttt{{[}192,12,15,26{]}} \\
\end{longtable}
}

\endgroup

Spatial compression consists of the 4× entry compression and three 2×
spatial merging steps:

\[
4\times2\times2\times2=32.
\]

The temporal lengths are:

\[
45\rightarrow23\rightarrow12.
\]

The output is split into two groups along the channel dimension, each
with shape \texttt{{[}96,12,15,26{]}}. Only the first group is taken as
\textbf{\(\mu\)}. The mean \(\mu\) is then normalized using the
per-channel means and standard deviations stored in the model.

\Needspace{6\baselineskip}

\subsection{5.4 Decoder: Reconstructing the Video Stage by
Stage}\label{flux3-decoder}

First undo the latent's per-channel normalization, then reconstruct the
RGB video through Linear projections, NA Blocks, and spatiotemporal
expansion. Shapes are written as \texttt{{[}C,T,H,W{]}}, omitting the
batch dimension.

\begingroup\NotesTableStyle\fontsize{10pt}{12pt}\selectfont\renewcommand{\arraystretch}{1.02}\setlength{\extrarowheight}{1pt}

{\def\LTcaptype{none} % do not increment counter
\begin{longtable}[]{@{}
  >{\raggedright\arraybackslash}p{(\linewidth - 4\tabcolsep) * \real{0.1700}}
  >{\raggedright\arraybackslash}p{(\linewidth - 4\tabcolsep) * \real{0.5500}}
  >{\raggedright\arraybackslash}p{(\linewidth - 4\tabcolsep) * \real{0.2800}}@{}}
\toprule\noalign{}
\begin{minipage}[b]{\linewidth}\raggedright
Step
\end{minipage} & \begin{minipage}[b]{\linewidth}\raggedright
Module and operation
\end{minipage} & \begin{minipage}[b]{\linewidth}\raggedright
Output shape
\end{minipage} \\
\midrule\noalign{}
\endhead
\bottomrule\noalign{}
\endlastfoot
Input & Latent after undoing normalization &
\texttt{{[}96,21,15,26{]}} \\
Entry & Position-wise Linear, \textbf{96\(\rightarrow\)2048} &
\texttt{{[}2048,21,15,26{]}} \\
Stage 1 temporal expansion & NA Block ×8 \(\rightarrow\)
\textbf{temporal expansion, 21\(\rightarrow\)41} \(\rightarrow\) NA
Block ×1 & \texttt{{[}2048,41,15,26{]}} \\
Stage 1 spatial expansion & \textbf{Spatial expansion,
2048\(\rightarrow\)1024} & \texttt{{[}1024,41,30,52{]}} \\
Stage 2 temporal expansion & NA Block ×8 \(\rightarrow\)
\textbf{temporal expansion, 41\(\rightarrow\)81} \(\rightarrow\) NA
Block ×1 & \texttt{{[}1024,81,30,52{]}} \\
Stage 2 spatial expansion & \textbf{Spatial expansion,
1024\(\rightarrow\)512} & \texttt{{[}512,81,60,104{]}} \\
Stage 3 & NA Block ×4 \(\rightarrow\) \textbf{spatial expansion,
512\(\rightarrow\)256} & \texttt{{[}256,81,120,208{]}} \\
Stage 4 & NA Block ×1, retaining 256 channels &
\texttt{{[}256,81,120,208{]}} \\
Exit & Position-wise Linear, \textbf{256\(\rightarrow\)48} &
\texttt{{[}48,81,120,208{]}} \\
Unpatchify & Rearrange \textbf{48=3×4×4} channels into RGB pixel blocks
& \texttt{{[}3,81,480,832{]}} \\
\end{longtable}
}

\endgroup

Each temporal expansion follows \textbf{\(T\rightarrow2T-1\)}, giving:

\[
21\rightarrow41\rightarrow81.
\]

\clearpage

\subsection{5.5 Downsampling, Upsampling, and First-Frame
Alignment}\label{flux3-sampling}

\begingroup\NotesTableStyle\fontsize{10pt}{12pt}\selectfont\renewcommand{\arraystretch}{1.02}\setlength{\extrarowheight}{1pt}

{\def\LTcaptype{none} % do not increment counter
\begin{longtable}[]{@{}
  >{\raggedright\arraybackslash}p{(\linewidth - 4\tabcolsep) * \real{0.2000}}
  >{\raggedright\arraybackslash}p{(\linewidth - 4\tabcolsep) * \real{0.5000}}
  >{\raggedright\arraybackslash}p{(\linewidth - 4\tabcolsep) * \real{0.3000}}@{}}
\toprule\noalign{}
\begin{minipage}[b]{\linewidth}\raggedright
Operation
\end{minipage} & \begin{minipage}[b]{\linewidth}\raggedright
Implementation
\end{minipage} & \begin{minipage}[b]{\linewidth}\raggedright
Shape change
\end{minipage} \\
\midrule\noalign{}
\endhead
\bottomrule\noalign{}
\endlastfoot
\textbf{Entry Patch Embedding} & \texttt{1×4×4\ Conv3d}, stride
\texttt{1×4×4}, \textbf{3\(\rightarrow\)256} & Divide \(H,W\) by 4;
\(T\) unchanged \\
\textbf{Spatial merging} & Concatenate \texttt{2×2} neighboring
positions: \texttt{C→4C} \(\rightarrow\) LayerNorm \(\rightarrow\)
Linear, \textbf{4C\(\rightarrow\)2C} & Halve \(H,W\); double \(C\) \\
\textbf{Temporal merging} & For odd lengths, repeat the first position
at the beginning \(\rightarrow\) concatenate pairs of positions
\(\rightarrow\) LayerNorm \(\rightarrow\) Linear,
\textbf{2C\(\rightarrow\)C} \(\rightarrow\) add a shortcut containing
the mean of the two positions & \(T\rightarrow\lceil T/2\rceil\);
\(C,H,W\) unchanged \\
\textbf{Spatial expansion} & LayerNorm \(\rightarrow\) Linear,
\textbf{C\(\rightarrow\)4C\_out} \(\rightarrow\) rearrange into
\texttt{2×2} spatial positions, where \texttt{C\_out=C/2} & Double
\(H,W\); halve \(C\) \\
\textbf{Temporal expansion} & LayerNorm \(\rightarrow\) Linear,
\textbf{C\(\rightarrow\)2C} \(\rightarrow\) add a shortcut consisting of
two copies of the input along the channel dimension \(\rightarrow\)
rearrange into the time dimension \(\rightarrow\) remove the first
position & \(T\rightarrow2T-1\); \(C,H,W\) unchanged \\
\textbf{Spatial Unpatchify} & Rearrange the 48 channels at each position
into a \texttt{3×4×4} pixel block & Multiply \(H,W\) by 4; output 3
channels \\
\end{longtable}
}

\endgroup

\textbf{Execution order:}

\begin{itemize}
\tightlist
\item
  Encoder Stage 3: spatial merging \(\rightarrow\) additional NA Block
  \(\rightarrow\) temporal merging. Stage 4 applies an additional NA
  Block \(\rightarrow\) temporal merging after its backbone.
\item
  The first two Decoder stages: temporal expansion \(\rightarrow\)
  additional NA Block \(\rightarrow\) spatial expansion.
\end{itemize}

\textbf{First-frame alignment:} encoding pads odd lengths by repeating
the first position, while decoding restores alignment by removing the
first position after expansion. Single-frame inputs also pass through
these modules, but their length remains 1:

\begin{verbatim}
Encoder: 1 -> repeat to 2 -> merge to 1
Decoder: 1 -> expand to 2 -> remove the first position, leaving 1
\end{verbatim}

After two temporal expansion steps, the output length is:

\[
T_{\mathrm{out}}=2(2T_z-1)-1=4(T_z-1)+1.
\]

\end{document}
